\section{Reproducibility Details}
\label{app:reproducibility}
\subsection{Initial Conditions}
\texttt{scripts/make\_decks.py} defines the IC distributions. For each sample index $r\in\{0,\ldots,99\}$ it draws a fixed two-channel random array with NumPy's PCG64 generator and seed sequence $[20261118,r]$. Each IC family maps the same random array to its own parameters, and the two fields use separate rows. Parameters are rounded to eight decimal places before the field is evaluated, so the saved numbers define the actual initial field.

Let $\alpha\sim\mathcal U[0.2,0.8]$, $m_x,m_y\sim\mathcal U[0.25,0.75]$, and $\sigma\sim\mathcal U[0.08,0.18]$, and define
\[
 B(x,y;m_x,m_y,\sigma)=
 \exp\!\left[-\frac{(x-m_x)^2+(y-m_y)^2}{2\sigma^2}\right].
\]
Each field uses one of the following:
\begin{align*}
 \text{Single bump:}\quad &\alpha w(x,y)B(x,y;m_x,m_y,\sigma),\\
 \text{Oscillations:}\quad &\frac{\alpha w(x,y)}{1.5}
 \big[1+0.5\cos(2\pi k_xx+\phi_x)\cos(2\pi k_yy+\phi_y)\big],\\
 \text{Three patches:}\quad &\alpha w(x,y)
 \frac{\sum_{\ell=1}^3 b_\ell B(x,y;m_{x,\ell},m_{y,\ell},\sigma_\ell)}
 {\sum_{\ell=1}^3 b_\ell}.
\end{align*}
Here $k_x,k_y$ are uniform on $\{1,2,3\}$, the phases are uniform on $[0,2\pi)$, and the patch weights $b_\ell$ are uniform on $[0.2,1]$. Patch centers and widths use the single-bump ranges. Every initial field lies in $[0,0.8]$. The same parameter bank is used for all equations and BCs. Linking different IC families through a shared seed is a design choice; it is not an optimal coupling between their distributions.

\subsection{Descriptions and Records}
The release contains 54 JSON problem files, the 300-sample IC parameter bank, and all 5,400 text and JSON descriptions. Running the generator with \texttt{--all-realizations} rebuilds the full set; without it, the generator writes three examples.

The text given to a model contains the two equations with all coefficients, the two explicit initial fields, the boundary data, the domain, and the time interval. The JSON version adds the IC family, sample index, grid, saved times, and observed frames. Solver settings, checksums, and diagnostics are stored in \texttt{benchmark/results/simulation\_manifest.json}. Model versions, preprocessing settings, token counts, and timings are stored next to the embeddings. Keeping these records out of the model text means that no solver setting or label can leak into a model input.

Below is the complete text given to a model for one sample of advection--diffusion with a single-bump IC and zero Neumann BC (long lines are wrapped for display). Its equivalent form would be $u_t=0.01(u_{xx}+u_{yy})-0.25u_x-0.10u_y$, and likewise for $v$. A model never has to guess what a named solver's defaults would be.
\begin{lstlisting}[basicstyle=\ttfamily\scriptsize,breaklines=true,columns=fullflexible,frame=single]
Equation:
\partial_t u=D_u\Delta u-c_x\partial_xu-c_y\partial_yu,\qquad \partial_t v=D_v\Delta v-c_x\partial_xv-c_y\partial_yv.

Parameters:
D_u=0.01000000,\quad D_v=0.00500000,\quad c_x=0.25000000,\quad c_y=0.10000000.

Initial condition:
\eta(s)=\begin{cases}\exp\!\left(4-\frac{1}{s(1-s)}\right),&0<s<1,\\0,&s\in\{0,1\}.\end{cases}
u(0,x,y)=0.41581429\,\eta(x)\eta(y)\,\exp\!\left[-\frac{(x-0.52007856)^2+(y-0.44353150)^2}{2(0.08503441)^2}\right].
v(0,x,y)=0.62013414\,\eta(x)\eta(y)\,\exp\!\left[-\frac{(x-0.74287585)^2+(y-0.27895931)^2}{2(0.13975413)^2}\right].

Boundary:
\partial_{\mathbf n}U(t,x,y)=0\quad\text{for }(x,y)\in\partial\Omega.

Domain:
\Omega=[0,1]^2,\qquad U=(u,v)^\top,\qquad \Delta=\partial_{xx}+\partial_{yy}.\quad x_p=(p+\tfrac12)/64,\ y_q=(q+\tfrac12)/64,\ p,q=0,\ldots,63.

Time:
0\leq t\leq1,\qquad t_n=n/99,\quad n=0,\ldots,99.\quad \text{Observed frames: }n\in\{0,25,50,74,99\}.
\end{lstlisting}

\subsection{Solver Settings}
Cell centers are $x_p=(p+1/2)/64$ and $y_q=(q+1/2)/64$ for $p,q=0,\ldots,63$, and arrays are ordered (field, $x$, $y$). The exact linear solutions in \cref{app:linear-reference} use the same stencil and ghost values as the solver, so they test the same grid equations. One solver serves all six equations, so no equation is tied to a different grid or boundary treatment.

Output times and solver steps differ even for diffusion. With $h=1/64$ and $\nu_{\max}=0.01$, an explicit five-point scheme would need $\delta t\le h^2/(4\nu_{\max})\approx0.00610$, while outputs are $1/99\approx0.01010$ apart. ETDRK4 handles diffusion exactly and uses $\delta t=1/396$ for the remaining terms. No trajectory failed, was dropped, or was clipped; stored values lie in $[\simMin,\ \simMax]$.

APEBench's periodic interface uses \texttt{num\_train\_samples=100} and \texttt{train\_temporal\_horizon=99} for 100 samples with 100 saved states \citep{apebench_base_docs}. Its default time steps and IC generators differ from ours, so the periodic check in \cref{app:validation} calls the Exponax solvers directly with our initial arrays.

\subsection{Embedding Protocol}
\label{app:embedding-protocol}
The model versions are \texttt{google/embeddinggemma-2} (revision \texttt{f9972ac}), \texttt{Qwen/Qwen3-VL-Embedding-2B} (\texttt{9f2f7e7}), and \texttt{Alibaba-NLP/gme-Qwen2-VL-2B-Instruct} (\texttt{9cfa641}); the full hashes are in the protocol file. Their model cards describe combined text and image inputs \citep{embeddinggemma2_modelcard,qwen3vl_embedding_modelcard,gme_modelcard}, and the Qwen3-VL and GME reports describe their training and evaluation \citep{li2026qwen3vl,zhang2024gme}.

Each model gets its documented input format, with one fixed prompt for every input. EmbeddingGemma 2 receives the symmetric prefix ``\texttt{task: sentence similarity | query: }'' on inputs with text, no prefix on images, and, for joint inputs, the description followed by the \texttt{<|image|>} placeholder; we use the checkpoint's mean pooling. Qwen3-VL-Embedding follows its official script: the instruction ``Represent the user's input.'', image before text, the generation prompt, last-token pooling, and image resizing by \texttt{qwen-vl-utils}. GME follows its reference code \texttt{gme\_inference.py}: the system prompt ``You are a helpful assistant.'', image before text, a final \texttt{<|endoftext|>} token, sequential positions, 256--1280 image tokens, and a limit of 1,800 tokens. We do not use separate query and document prompts, because our score in \cref{eq:embedding-score} treats both sides of a pair the same way. All models run in bfloat16 with batch size one, which this CPU supports natively, and every output is renormalized in float32 at full size. No input was truncated (\cref{tab:model-inputs}).

Before encoding any test input, we checked each interface on a pilot image; the results are stored in \texttt{benchmark/embedding\_interface\_checks.json}. Our GME code produces the same tokens and image grid as the official code run under transformers 4.51.3, and the embeddings agree to a cosine similarity of at least 0.997; the remaining difference comes from bfloat16 arithmetic in different library versions. Our Qwen3-VL code agrees with the sentence-transformers interface on the model card (cosine at least 0.998). Putting the text before the image instead changes the Qwen3-VL joint embedding a lot (cosine 0.905 with the official order), so input order is part of the protocol.

\subsection{What Was Encoded}
On a CPU with two cores and no GPU, encoding all $3\times3\times5400=48{,}600$ main inputs was not possible. The protocol therefore fixed a smaller set in advance: samples 0--9 of all 54 problems with all three input types (1,620 inputs per model), and samples 0--2 for the controls (162 reversed images, 162 rewritten descriptions, 162 mismatched pairs, 18 initial-frame images, and four heat-equation descriptions per model). Every model and baseline is compared with the reference computed on the same samples. The full run only needs \texttt{scripts/run\_embeddings.py --realizations 100}.

\subsection{Statistics and Release}
Samples with index 100 or higher were used only for setup and image scaling; the test set is samples 0--99. Bootstrap intervals resample the encoded samples (2,000 times) and are rough with ten samples. Every encoding attempt is logged; a failure would appear in the embedding records, and none occurred.

The release contains the problem files, the parameter bank, all descriptions, the solver and its checks, the pilot and run records, the per-sample reference terms, all embeddings, the analysis code, and the generated tables. It also includes the 540 images that were encoded in the main runs (samples 0--9). The five-frame arrays (0.9\,GB) and the remaining images are regenerated by \texttt{scripts/run\_simulations.py} and checked against the recorded hashes.
